% 13_negative
\section{无收益与已回退的尝试}

本节汇总在真机上测平或测负、以及实现后回退的尝试，覆盖调度、内存、启动、频率与 NPU 五处路径。每一项尝试都在开发板上完成 A/B 对照，通过者予以保留，测平或测负者如实记录于此。明确记录这些尝试，为当前平台划定了边界，也避免日后重复无效路径。多数尝试在更换硬件类型或补足某项前置条件后仍有望转为正收益，其适用边界随条目一并列出。

下表按调度与内存、启动与频率、NPU 三组，列出各项尝试的现象、失效原因与适用边界。

\begin{table}[H]\centering\small
\begin{tabular}{@{}>{\raggedright\arraybackslash}p{2.5cm}>{\raggedright\arraybackslash}p{2.6cm}>{\raggedright\arraybackslash}p{5.6cm}>{\raggedright\arraybackslash}p{3.4cm}@{}}
\toprule
尝试 & 现象 & 为何无收益或回退 & 适用边界 \\
\midrule
全量 work-stealing 负载均衡器（idle-pull、push 与 wake-redirect） &
单线程吞吐减半；八线程饱和时与关闭时基本持平 &
七个空闲核心上，三条迁移路径反复搬动唯一的工作线程，无核心能维持缓存热度；大小核架构上，分离一对热任务与令独立任务保持缓存热度两个目标直接冲突。已回退，运行期迁移仅以默认关闭的 feature 保留 &
核心众多、任务长期饱和、迁移成本相对计算量可忽略的同构机器，或补足按缓存热度与容量设定的迁移门控之后 \\
\addlinespace
自适应 haltpoll &
唤醒延迟未能稳定优于固定时长轮询 &
停机时长信号受 GIC-SGI 与 WFI 交互的异常特性干扰，判据失真。已回退为固定时长自旋（idle-poll）后再进入 WFI；idle-poll 现随并发工作线程的唤醒铺开在出厂放置配置中启用 &
平台停机时长信号可信、且无 SGI-WFI 干扰时 \\
\addlinespace
按需缺页拷贝（去掉 fork 前的 \texttt{populate\_area}） &
fork 密集的 \texttt{-P} 负载显著变慢，线程路径仅获小幅收益 &
\texttt{-T} 线程路径的瓶颈源自并发度，持锁时长与缺页模型均不构成主因。已回退 &
缺页并发不构成瓶颈、且惰性建立映射能显著节省未触碰页开销的场景 \\
\addlinespace
early-MMU 启动重排（提前开启 MMU 与缓存） &
kernel 到 shell 用时与未改动基本相同 &
启动早期的未缓存窗口在启动低频下受限于 CPU 计算，对计算密集工作启用缓存不省时间。已回退，保留启动计时诊断 &
该窗口由内存延迟主导，或启动核心已运行于高频时 \\
\addlinespace
探测让步延时（yield 式探测等待） &
并发设备探测耗时增加，整机启动更慢 &
tick 粒度将众多小等待分别向上取整至一个 tick，逐条放大每条并发探测。已回退 &
单次等待远长于一个 tick，或改用高分辨率定时器进行让步时 \\
\addlinespace
AVS/PVTM 顶频抬升 &
A76 顶档上限无法提升 &
顶部 A76 OPP 全部锁定于固定电压，无 AVS 分箱，主线亦无对应的 Rockchip AVS 驱动；真正提升顶频的机制是 PVTPLL 环长切换，A76 治理上限 \nDvfsAseven~MHz（保留待温控可至 \nDvfsAsevenTop~MHz）、A55 \nDvfsAfive~MHz，逐核约达 Linux 的 \nPerCoreRatio，见频率电压一节。已回退 &
芯片提供非空 AVS 分箱、且存在可读取分箱电压的驱动时 \\
\addlinespace
NPU 张量零拷贝（把张量重排移到用户态） &
端到端推理时间反而上升 &
单核推理下仅改变重排的执行位置，重排开销未减少，且新增一次搬运。已回退 &
重排能真正并行至另一核心，或交由硬件承担时 \\
\addlinespace
第三 NPU 核 &
由两核启用至三核推理无收益 &
当前模型负载约占用一个核心，增配核心无法充分占用。已保持两核 &
模型足够大、单核算力成为瓶颈时 \\
\addlinespace
NPU 定频（锁住 governor） &
无变化 &
NPU 已稳定运行于其额定工作频率，锁频未改变实际工作点 &
NPU 频率原本会被 ondemand 治理压低或产生抖动时 \\
\bottomrule
\end{tabular}
\caption{在真机上测平或测负、以及实现后回退的尝试：现象、失效原因与适用边界}
\end{table}

几条回退指向同一批平台事实。大小核架构上，分离一对热任务与令独立任务保持缓存热度两个目标天然冲突，全量迁移在七个空闲核心的场景下只反复搬动唯一的工作线程，这也是\ours{占用感知的放置}仅在 spawn 与 wake 两处执行、运行期迁移默认关闭的原因。自适应 haltpoll 的不稳定与朴素唤醒铺开的高代价同源，均受制于该 SoC 毫秒量级的 reschedule SGI 唤醒 WFI 核心的异常特性；改以固定时长自旋的 idle-poll 抹平这一异常之后，唤醒铺开转为收益并随之出厂，粗粒度 \texttt{wake\_affine} 仍把生产者消费者的\keyres{唤醒延迟收在 \nWake~µs（Linux \nWakeLin~µs）}，铺开机制见调度一节。NPU 三项无收益同源于单核已能容纳当前模型负载，改变重排位置、增配核心或锁频都不触及实际工作点。AVS 顶频属于固件与硬件的外部约束，真正提升顶频的机制是 PVTPLL 环长切换，见频率电压一节。这些负结果不改变整体判断，缩小了后续可着手的空间，也标出更换平台类型时应当重试的候选。
